\chapter{Métaux}\label{ch:m3:metals}

Au petit matin du mardi 8 mars 2022, le contrat à trois mois sur le nickel du London
Metal Exchange est passé de juste en dessous de 50\,000 dollars la tonne à 101\,365 dollars à
06:08, soit plus que doublé en cinq heures environ~; la veille, il avait ouvert juste en dessous
de 30\,000. À 08:15, la bourse a suspendu les négociations. À 12:05, elle a annulé chaque transaction sur le nickel
conclue depuis minuit. Sa chambre de compensation avait calculé que, si les transactions avaient été maintenues, les membres auraient dû
déposer quelque 19.75 milliards de dollars de marge à très bref délai, avec un risque de
défauts multiples. Les métaux de base se négocient sur une bourse dont les contrats, les entrepôts et les règles
diffèrent de ceux de tout autre marché à terme, et la semaine du short squeeze sur le nickel
les a tous mis à l'épreuve. Ce chapitre couvre les \omterm{def:m3:metals:prompt}{prompt dates} et les prix du marché londonien des métaux, son système
d'entrepôts, l'épisode du nickel, les deux marchés de l'or, à Londres et à New York, et le minerai de fer.

\section{Le marché londonien des métaux~: prompt dates et prix}

Le London Metal Exchange négocie l'aluminium, le cuivre, le zinc, le nickel, le plomb, l'étain et d'autres
métaux. Ses contrats ne sont pas des contrats à terme mensuels~: ce sont des contrats de livraison à une date
parmi un ensemble dense de dates, une structure héritée du temps où il fallait aux navires
pour acheminer le métal jusqu'à Londres.

\begin{definition}[Prompt date, prix à trois mois]\label{def:m3:metals:prompt}
Une \emph{prompt date}\index{prompt date} est une date à laquelle un contrat du LME peut être réglé par
livraison de warrants contre paiement. Il y en a un pour chaque jour ouvré, de demain
(tom) à trois mois, un pour chaque mercredi jusqu'à six mois, et un le
troisième mercredi de chaque mois au-delà, dans une limite qui dépend du métal. Le
\emph{prix à trois mois}\index{prix à trois mois} (3M) est le prix de la date prompt située à peu près
à trois mois de la date de transaction~; c'est le contrat du LME le plus liquide et la référence pour
la plupart des opérations de couverture.
\end{definition}

\begin{omfigure}
\begin{tikzpicture}
\begin{axis}[width=11cm,height=3.6cm,
  xlabel={mois après la date de transaction},
  tick label style={font=\small},label style={font=\small},
  xmin=0,xmax=27,ymin=0.5,ymax=3.5,ytick={1,2,3},yticklabels={mensuel,hebdomadaire,quotidien},
  grid=major,grid style={gray!25},
  table/col sep=comma]
\addplot[omDef,only marks,mark=|,mark size=3pt] table[x=months,y=row]{figdata/markets-3/08-metals/prompts.csv};
\end{axis}
\end{tikzpicture}

\omcaption{\omterm{def:m3:metals:prompt}{Prompt dates} disponibles le 24 septembre 2026 selon la structure de la
\cref{def:m3:metals:prompt}, sur 27 mois~: 65 dates quotidiennes, 14 hebdomadaires et 21 mensuelles,
les week-ends uniquement (sans calendrier de jours fériés). Données~: la construction du chapitre.}\label{fig:m3:metals:prompts}
\end{omfigure}

Une structure de dates dense permet à un consommateur d'ajuster sa couverture au jour où il a besoin du métal, mais elle
disperse la liquidité. Les traders déplacent leurs positions d'une date à l'autre avec des carries, le nom
que donne le marché des métaux à un calendar spread (livre~1, chapitre~19), et le plus court d'entre eux est le
tom-next du livre~2, chapitre~16~: vendre pour demain et racheter le jour suivant
revient à prêter du métal pour une journée.

\begin{example}[Un carry assimilé à un taux d'intérêt]\label{ex:m3:metals:carry}
Le cuivre au comptant est à 9\,800 dollars la tonne et à trois mois à 9\,880~: un contango de 80 dollars, soit 0.82~\% du
cours au comptant pendant environ 91 jours, soit 3.3~\% par an. Un détenteur de métal qui vend au comptant et achète à trois
mois (prête du métal) renonce à 80 dollars la tonne et économise le coût de financement et de stockage du
métal pendant trois mois~; si l'argent et le stockage coûtent plus de 3.3~\% par an, le prêt est le meilleur
usage du métal. Une backwardation de même ampleur rapporterait au prêteur 80 dollars à la place~: le
marché rémunère le métal dès aujourd'hui.
\end{example}

\begin{definition}[Prix officiel du LME]\label{def:m3:metals:official}
Le \emph{prix officiel du LME}\index{prix officiel du LME} d'un métal est le prix fixé chaque
jour ouvré à partir des derniers cours acheteurs et vendeurs de la seconde séance du
parquet à la criée de la bourse (le Ring), pour le comptant, à trois mois et certaines dates ultérieures~; les contrats physiques du secteur
se fixent par rapport à lui, comme les cargaisons de pétrole par rapport à l'évaluation d'une agence.
\end{definition}

\section{Entrepôts, warrants et files d'attente}

Les contrats du LME se livrent par transfert de propriété du métal stocké dans des entrepôts que
la bourse a approuvés, sous forme de documents qui circulent comme des titres de propriété.

\begin{definition}[Warrant d'entrepôt, warrant annulé, file de sortie]\label{def:m3:metals:warrant}
Un \emph{warrant d'entrepôt}\index{warrant d'entrepôt} est un titre de propriété sur un lot précis de
métal d'une marque et d'une qualité spécifiées, stocké dans un entrepôt agréé par le LME, et le moyen
de livraison des contrats du LME. Un \emph{warrant annulé}\index{warrant annulé} est un warrant dont le propriétaire a
averti l'entrepôt qu'il retirera le métal~: le lot quitte le système du LME et attend
d'être chargé pour la sortie. La \emph{file de sortie}\index{file de sortie} est l'arriéré de métal annulé
en attente de quitter un entrepôt lorsque les annulations dépassent sa capacité quotidienne de chargement.
\end{definition}

\begin{omfigure}
\centering
\begin{tikzpicture}[font=\small,
  st/.style={draw,thick,rounded corners=2pt,align=center,minimum height=1cm,text width=2.4cm},
  arr/.style={-{Stealth[length=2.2mm]},thick}]
\node[st,draw=omExo,fill=omExo!8] (in) at (0,0) {métal livré\\dans un entrepôt};
\node[st,draw=omDef,fill=omDef!8] (on) at (3.4,0) {sous warrant~:\\livrable};
\node[st,draw=omProp,fill=omProp!8] (can) at (6.8,0) {warrant\\annulé};
\node[st,draw=omThm,fill=omThm!8] (q) at (10.2,0) {file de\\sortie};
\draw[arr] (in) -- (on);
\draw[arr] (on) -- (can);
\draw[arr] (can) -- (q);
\draw[arr] (q.south) -- ++(0,-0.7) node[below,font=\footnotesize] {chargé vers un consommateur};
\draw[arr] (on.south) -- ++(0,-0.7) node[below,font=\footnotesize] {négocié au LME};
\end{tikzpicture}

\omcaption{Le parcours d'un lot de métal dans le système du LME. Seul le métal sous warrant peut être
livré contre un contrat~; le métal annulé attend de sortir, et une longue file immobilise du métal
que le marché compte comme stock. Schéma de principe.}\label{fig:m3:metals:warrants}
\end{omfigure}

Les stocks sous warrant sont ce que les vendeurs à découvert peuvent livrer~; les stocks annulés sont du métal déjà réclamé. Une
baisse des stocks sous warrant, ou une forte hausse des annulations, tend les dates proches et peut faire basculer la
courbe en backwardation. Les longues files ont aussi servi à accroître les loyers perçus et les primes
que les consommateurs paient pour le métal prompt~; le règlement de la bourse plafonne désormais le loyer que les entrepôts peuvent facturer pour
le métal placé dans de longues files.

\begin{definition}[règle de prêt du LME]\label{def:m3:metals:lending}
Une \emph{règle de prêt du LME}\index{règle de prêt du LME} impose au détenteur d'une position longue dominante
(de warrants et de contrats à échéance proche) de prêter (vendre à la date la plus proche et racheter à la
date la plus lointaine) à une backwardation ne dépassant pas un niveau fixé, de sorte que détenir la majeure partie
du métal livrable ne puisse pas servir à coincer les vendeurs à découvert à volonté.
\end{definition}

\section{Nickel, mars 2022}

\begin{omfigure}
\begin{tikzpicture}
\begin{axis}[width=11cm,height=5.2cm,
  ylabel={milliers de dollars US par tonne},
  tick label style={font=\small},label style={font=\small},
  xmin=0.5,xmax=8.5,ymin=0,ymax=115,grid=major,grid style={gray!25},
  xtick={1,2,3,4,5,6,7,8},
  xticklabels={ven. ouverture,ven. clôture,lun. ouverture,lun. clôture,mar. 04:49,mar. 06:08,mar. 07:00,mar. 08:15},
  xticklabel style={rotate=35,anchor=north east,font=\footnotesize},
  table/col sep=comma]
\addplot[omThm,very thick,mark=*,mark size=1.6pt] table[x=i,y=price]{figdata/markets-3/08-metals/nickel.csv};
\node[font=\footnotesize,anchor=east] at (axis cs:5.8,101.4) {pic 101\,365};
\node[font=\footnotesize,anchor=north west] at (axis cs:4.1,46) {48\,078};
\node[font=\footnotesize,anchor=south west] at (axis cs:0.8,36) {27\,080 à 28\,919};
\node[font=\footnotesize,anchor=north,align=center] at (axis cs:7.5,72) {suspendu\\à 08:15};
\end{axis}
\end{tikzpicture}

\omcaption{Nickel à trois mois au LME, du 4 au 8 mars 2022~: les cours cités dans l'arrêt de la Court of Appeal,
présentés dans l'ordre de négociation (week-end omis~; les niveaux après 07:00 repris de la mention «~
constamment au-dessus de 80\,000~»). Toutes les transactions conclues après minuit le 8 mars ont été
annulées. Source~: \emph{R (Elliott Associates) v London Metal Exchange} [2024] EWCA Civ
1168.}\label{fig:m3:metals:nickel}
\end{omfigure}

Les cours ont progressé en trois temps (\cref{fig:m3:metals:nickel}). Le vendredi 4 mars, le nickel à trois
mois a progressé de 6.8~\%, de 27\,080 à 28\,919 dollars, après l'invasion de l'Ukraine par la Russie, une
semaine plus tôt~; la chambre de compensation a appelé environ 2.6~milliards de dollars de marge intrajournalière, soit
40~\% au-dessus de son précédent record. Le lundi 7 mars, le cours a grimpé de 66~\% de clôture à clôture, jusqu'à 48\,078,
une variation quotidienne près de cinq fois supérieure à toute autre des vingt années précédentes, et les appels de marge ont totalisé environ 7.05
milliards de dollars~; trois membres compensateurs n'ont pas payé à temps. Dans les premières heures du mardi, il
a grimpé jusqu'à 101\,365. À 04:49, le cours étant proche de 60\,000, le personnel de la bourse a suspendu les
limites de variation de cours, car des transactions se concluaient au-dessus de celles-ci~; le cours est resté au-dessus de 80\,000 à partir
de 07:00 jusqu'à la suspension.

Selon la cour, la cause était un short squeeze (livre~1, chapitre~16)~: de fortes positions courtes avaient été
constituées par plusieurs intervenants, dont un producteur chinois de nickel,
Tsingshan, sur le marché de gré à gré~; à mesure que le cours montait, les vendeurs à découvert ont dû racheter pour faire face aux appels de marge,
et d'autres ont acheté devant eux. Le marché a rouvert le 16 mars, après que les banques du producteur
ont mis en place un soutien. L'annulation a été contestée en justice par des sociétés de trading qui avaient
acheté pendant la nuit~; la High Court, puis la Court of Appeal (2024), ont jugé que la
bourse avait agi légalement.

\begin{dated}{2026-09}{Le constat du régulateur}\label{dat:m3:metals:fca}
Le 19 mars 2025, la Financial Conduct Authority a infligé au London Metal Exchange
une amende de 9\,245\,900~livres sterling, sa première action répressive contre une bourse d'investissement reconnue,
pour ne pas avoir assuré un marché ordonné en mars 2022. Elle a constaté que seul du personnel junior, non formé
à repérer toutes les causes de désordre, était de permanence pendant les heures asiatiques, et que ce personnel
s'efforçait d'accompagner les hausses, notamment en désactivant les limites de variation de cours. La bourse a transigé et a obtenu une
remise de 30~\%.
\end{dated}

\begin{proposition}[Marge sur une position courte coincée]\label{prop:m3:metals:margin}
Une position courte de $q$ tonnes, réévaluée de $P_0$ à $P_1$, doit une marge de variation $q\,(P_1 - P_0)$ en
espèces le jour même. Une position courte de 10\,000 tonnes a dû verser 191.59 millions de dollars le lundi, et
aurait dû verser 532.87 millions de dollars de plus au pic du mardi~: 724.46 millions de dollars sur une position
dont la valeur était de 289 millions de dollars à la clôture du vendredi.
\end{proposition}

\begin{proof}
$10\,000 \times (48\,078 - 28\,919)$ et $10\,000 \times (101\,365 - 48\,078)$~; la position valait
$10\,000 \times 28\,919$ à la clôture du vendredi.
\end{proof}

Un producteur qui a vendu à terme sa production future est économiquement couvert~: le métal
qu'il produira prend de la valeur avec sa position courte. Sa marge, en revanche, se paie en espèces dès maintenant, et la
plus-value du métal n'arrive qu'à la vente. Le short squeeze du nickel est l'exemple le plus net d'une couverture
saine en valeur et fatale en liquidité, un thème du \cref{ch:m3:when-markets-break-together}.

\section{Or~: Londres contre New York}

L'or se négocie en deux lieux aux conventions différentes. À Londres, il se traite de gré à gré,
avec règlement en \omterm{def:m3:metals:loco}{or non alloué}.

\begin{definition}[Loco London, or non alloué]\label{def:m3:metals:loco}
\emph{Loco London}\index{loco London} est la convention selon laquelle un cours concerne de l'or (ou de l'argent)
livré à Londres, réglé par le système de compensation londonien des comptes non alloués.
\emph{or non alloué}\index{or non alloué} est un solde d'or détenu sur un compte auprès d'une banque
de compensation~: le détenteur ne possède aucun lingot précis, seulement une créance sur la banque pour ce poids de métal, et
est un créancier chirographaire de la banque.
\end{definition}

\begin{definition}[LBMA Gold Price]\label{def:m3:metals:lbma}
Le \emph{LBMA Gold Price}\index{LBMA Gold Price} est l'indice de référence du cours au comptant de l'\omterm{def:m3:metals:loco}{or non alloué},
\omterm{def:m3:metals:loco}{loco London}, fixé deux fois par jour ouvré à Londres, lors d'enchères électroniques conduites par un
administrateur indépendant.
\end{definition}

\begin{dated}{2026-09}{Les enchères de Londres}\label{dat:m3:metals:lbma}
Les enchères \omterm{def:m3:metals:lbma}{LBMA Gold Price} se tiennent à 10:30 et 15:00, heure de Londres, et sont administrées par ICE Benchmark
Administration, qui les organise depuis mars 2015.
\end{dated}

À New York, l'or se négocie sous forme de contrats à terme sur le COMEX, livrés en lingots de 100 onces dans des chambres fortes agréées.
Les deux marchés sont reliés par l'exchange for physical (livre~1, chapitre~21)~: un trader acheteur de contrats à terme
et vendeur de swaps sur l'or de Londres les échange l'un contre l'autre, et d'ordinaire les contrats à terme se négocient à une
prime sur Londres égale au carry et à de petits frais, soit environ 1.50~\$ par once. En mars 2020, le lien
s'est rompu. Les raffineries suisses qui refondaient les lingots de 400 onces de Londres en barres plus petites
livrées à New York ont fermé, et les vols de passagers qui transportent l'or ont été cloués au sol~; les 25--26 mars, les contrats à terme d'avril
se sont négociés à plus de 70~\$ au-dessus de Londres, puis à 50~\$.
La bourse et l'association londonienne ont répondu par un contrat à terme livrable en lingots de 100 onces, de 400 onces
ou d'un kilogramme.

\begin{omfigure}
\centering
\begin{tikzpicture}[font=\small,
  mk/.style={draw,thick,rounded corners=3pt,align=center,minimum width=3.9cm,minimum height=1.6cm},
  arr/.style={{Stealth[length=2.2mm]}-{Stealth[length=2.2mm]},thick}]
\node[mk,draw=omDef,fill=omDef!8] (ldn) at (0,0) {Marché de gré à gré de Londres\\or non alloué, Loco London\\lingots de 400 onces};
\node[mk,draw=omThm,fill=omThm!8] (ny) at (8.4,0) {Future de New York\\chambres fortes du COMEX\\lingots de 100 onces};
\draw[arr] (ldn.east) -- node[above=3pt,font=\footnotesize] {EFP} node[below=3pt,font=\footnotesize,align=center] {env.\\1.50~\$} (ny.west);
\node[font=\footnotesize,align=center,anchor=north] at (4.2,-1.2) {lien physique~: les raffineurs refondent les lingots,\\et les vols de passagers les transportent};
\end{tikzpicture}

\omcaption{Les deux marchés de l'or et le lien qui les unit. Le spread EFP reste proche de son coût
tant que les lingots peuvent être refondus et transportés~; en mars 2020, les deux ont cédé en même temps. Schéma~; source du
niveau normal et de la désorganisation de 2020~: Singapore Bullion Market Association,
2020.}\label{fig:m3:metals:gold}
\end{omfigure}

\section{Minerai de fer}

Le minerai de fer se négocie selon deux types de contrats très différents. À Singapour, il se négocie sous forme de contrats à terme réglés en espèces et de swaps, indexés sur la moyenne mensuelle
de l'évaluation, par une \omterm{def:m3:physical-commodity-markets:pra}{agence de cotation des prix}, de fines à 62\% de teneur en fer livrées en Chine~; à Dalian,
il se négocie sous forme de contrats à terme à livraison physique, par lots de 100 tonnes, ouverts aux investisseurs étrangers depuis
mai 2018. Le spread entre les deux constitue un marché à part entière, avec un \omterm{def:m3:physical-commodity-markets:basisrisk}{risque de base} du type de celui
rencontré au \cref{ch:m3:physical-commodity-markets}.

\section{Tutoriel~: le calendrier des prompt dates et une position short coincée}\label{tut:m3:metals:prompts}

\begin{tutorial}
\textbf{Objectif.} Construire le calendrier des \omterm{def:m3:metals:prompt}{prompt dates} et calculer la marge d'une position short pendant le
short squeeze sur le nickel. \textbf{État final~:} la \cref{fig:m3:metals:prompts} et les marges de la
\cref{prop:m3:metals:margin}.
\begin{tutsteps}
\item \textbf{\omterm{def:m3:metals:prompt}{Prompt dates}.} Quotidiennes jusqu'à trois mois, hebdomadaires les mercredis jusqu'à six mois, puis les troisièmes
mercredis au-delà.
\omcode{code/firm/prompts/firm_prompts.py}{59}{79}{Les prompt dates d'une date de transaction.}\label{lst:m3:metals:prompts}
\item \textbf{Le short squeeze.} Marge de variation jour par jour sur les cours retenus par le jugement.
\omcode{code/markets-3/08-metals/python/m3_metals.py}{28}{34}{Marge de variation d'une position short pendant les 7 et 8 mars 2022.}\label{lst:m3:metals:squeeze}
\item \textbf{Exécuter} \texttt{m3\_metals.calendar()}, \texttt{short\_squeeze()} et \texttt{fig\_metals.py}.
\end{tutsteps}
\textbf{Pour aller plus loin.} Ajouter le calendrier des jours fériés de la bourse~; chiffrer un carry du comptant à trois
mois à partir de deux cotations et le convertir en taux d'intérêt implicite annualisé.
\end{tutorial}

\section{Construction~: le calendrier des prompt dates et l'estimateur de marge}\label{bld:m3:metals:prompts}

\begin{build}
\bfield{Rôle} La mini-entreprise couvre son exposition aux métaux au LME~: elle doit savoir quelles dates existent aujourd'hui,
ce que coûte un carry et de combien de liquidités ses positions auront besoin si les cours bougent.
\bfield{Interface} \texttt{prompt\_dates(trade, months, holidays)} renvoie les dates quotidiennes, hebdomadaires, mensuelles,
de comptant et à trois mois~; \texttt{carry(near, far)}~; \texttt{variation\_margin(position, old,
new)}~; les fonctions utilitaires du calendrier \texttt{add\_months}, \texttt{lme\_roll},
\texttt{three\_month\_date}, \texttt{third\_wednesday}.
\bfield{Règles} Les jours ouvrés excluent les week-ends et les jours fériés donnés~; une \omterm{def:m3:metals:prompt}{prompt date} tombant un jour non ouvré
est reportée au jour ouvré suivant, mais un samedi est reporté au vendredi, et une date à trois mois repoussée dans le quatrième mois
tombe le dernier jour ouvré du troisième
(Trading Regulations 8.4.1 et 8.4.2 du LME)~; les prompts mensuelles commencent après la dernière prompt hebdomadaire.
\bfield{Tests d'acceptation} \texttt{code/firm/prompts/tests/}~: arithmétique des mois en fin de mois~; les reports du samedi,
du dimanche et de fin de mois~; la
structure quotidienne, hebdomadaire et mensuelle~; le signe de la marge d'une position short.
\bfield{Pour aller plus loin} Charger le calendrier de négociation publié par la bourse~; marge sous un scénario de
plusieurs jours de variations avec appels de marge intrajournaliers~; limites de la règle de prêt sur une position longue dominante.
\end{build}

\begin{omsources}
\item Court of Appeal, \emph{R (Elliott Associates) v London Metal Exchange} [2024] EWCA Civ 1168.
\item Financial Conduct Authority, communiqué de presse du 20 mars 2025 et Final Notice à l'attention du London
Metal Exchange.
\item London Metal Exchange, prompt date structure~; position management and lending rules~; warehouse
policy~; \emph{Rules and Regulations} (1 juin 2026), Part 3, Trading Regulation 8, prompt dates (copie Internet Archive).
\item LBMA, le LBMA Gold Price et loco London~; ICE Benchmark Administration.
\item J.~East, ``The day the EFP broke'', Singapore Bullion Market Association, 26 mars 2020.
\item SGX iron ore contract rules~; Dalian Commodity Exchange, ouverture des contrats à terme sur minerai de fer
aux investisseurs étrangers (2018).
\end{omsources}

\section{Exercices}

\begin{exercise}[$\star$]\label{exo:m3:metals:1}
Un jeudi, quelle date est le tom et laquelle est le comptant, si aucun jour férié n'intervient~?
\end{exercise}

\begin{exercise}[$\star$]\label{exo:m3:metals:2}
Un consommateur prête 500 tonnes de cuivre en tom-next. Que fait-il à chacune des dates~?
\end{exercise}

\begin{exercise}[$\star$]\label{exo:m3:metals:3}
Pourquoi une hausse des \omterm{def:m3:metals:warrant}{warrants annulés} tend-elle à resserrer les dates proches~?
\end{exercise}

\begin{exercise}[$\star\star$]\label{exo:m3:metals:4}
Une position short de 2\,000 tonnes de nickel a été valorisée au cours de clôture du vendredi. Quelle marge de variation devait-elle
le lundi, et combien aurait-elle dû au plus haut de la séance du mardi~?
\end{exercise}

\begin{exercise}[$\star\star$]\label{exo:m3:metals:5}
De combien le nickel a-t-il monté de l'ouverture du vendredi au plus haut du mardi~? Donner la variation de chaque journée.
\end{exercise}

\begin{exercise}[$\star\star$]\label{exo:m3:metals:6}
En mars 2020, pourquoi un trader ne pouvait-il pas simplement acheter de l'or de Londres à 1\,600~\$ et le livrer contre
des contrats à terme du COMEX à 1\,670~\$~?
\end{exercise}

\begin{exercise}[$\star\star\star$]\label{exo:m3:metals:7}
\emph{Programmation.} Avec \texttt{prompt\_dates}, compter les prompts quotidiennes, hebdomadaires et mensuelles disponibles le
24 septembre 2026, jusqu'à 27 mois.
\end{exercise}

\begin{exercise}[$\star\star\star$]\label{exo:m3:metals:8}
\emph{Repérer l'erreur.} «~Le producteur était entièrement couvert, donc la flambée des prix ne pouvait pas lui nuire.~»
\end{exercise}

\section{Problème~: mardi 8 mars 2022}

\begin{problem}[{Problème du week-end --- un short dans la squeeze du nickel}]\label{pb:m3:metals:1}
Un trader sur métaux est short de 10\,000 tonnes de nickel à trois mois, en prévision de la production attendue. Utiliser les
prix de la \cref{fig:m3:metals:nickel}.

\medskip\noindent\textbf{Partie I --- La position.}
\begin{enumerate}
\item Quelle était la valeur de la position à la clôture du vendredi~?
\item Combien la hausse du vendredi a-t-elle coûté en appels de marge~?
\item Combien a coûté la hausse du lundi~?
\item Combien aurait coûté le sommet du mardi, si les transactions avaient tenu~?
\item Donner le besoin total de liquidités entre l'ouverture du vendredi et le sommet.
\end{enumerate}

\medskip\noindent\textbf{Partie II --- Le marché.}
\begin{enumerate}[resume]
\item Quelle a été l'ampleur des appels de marge de la chambre de compensation les 4 et 7 mars~?
\item Quels montants les membres auraient-ils dû déposer le 8 mars si les transactions avaient tenu~?
\item Pourquoi les limites de variation de cours ont-elles été suspendues à 04:49~?
\item Pourquoi la bourse a-t-elle annulé les transactions plutôt que de simplement les suspendre~?
\item Qui a perdu à cause de l'annulation, et qu'ont décidé les tribunaux~?
\end{enumerate}

\medskip\noindent\textbf{Partie III --- Le mécanisme.}
\begin{enumerate}[resume]
\item Pourquoi une position courte importante et connue incite-t-elle d'autres acteurs à acheter~?
\item Comment une règle de prêt limite-t-elle les squeezes sur la bourse, et pourquoi n'a-t-elle pas empêché celle-ci~?
\item Pourquoi la squeeze a-t-elle débuté sur le marché de gré à gré~?
\item Quel rôle ont joué les faibles volumes d'échanges pendant les heures asiatiques~?
\item Qu'est-ce qui a permis la réouverture du marché le 16 mars~?
\end{enumerate}

\medskip\noindent\textbf{Partie IV --- Jugement.}
\begin{enumerate}[resume]
\item Que le trader aurait-il dû détenir en plus de son short pour survivre à la semaine~?
\item Qu'auraient changé les conclusions du régulateur dans le fonctionnement de la bourse~?
\item Une bourse doit-elle jamais annuler des transactions qui n'étaient pas des erreurs~?
\item Énoncer le \emph{résultat nommé}~: la marge due par un short de 10\,000 tonnes, de la clôture du lundi
jusqu'au sommet, et au total depuis la clôture du vendredi.
\item En une phrase~: qu'est-ce qui met à mal un producteur couvert lors d'une squeeze~?
\end{enumerate}
\end{problem}

\section{Questions d'entretien}

\begin{interviewq}[$\star$ \iqroles{trader}]\label{iq:m3:metals:1}
En quoi un contrat du LME diffère-t-il d'un contrat à terme mensuel du CME~?
\end{interviewq}

\begin{interviewq}[$\star$ \iqroles{trader, risque}]\label{iq:m3:metals:2}
Qu'est-ce qu'une short squeeze, et qu'est-ce que l'épisode du nickel a ajouté à votre compréhension de ce phénomène~?
\end{interviewq}

\begin{interviewq}[$\star\star$ \iqroles{chercheur}]\label{iq:m3:metals:3}
Les stocks sous warrant chutent de moitié en un mois. Que devient la courbe, et comment la négocieriez-vous~?
\end{interviewq}

\begin{interviewq}[$\star\star$ \iqroles{risque}]\label{iq:m3:metals:4}
Concevez le test de stress de liquidité d'un producteur qui couvre sa production avec des contrats à terme de bourse.
\end{interviewq}

\begin{interviewq}[$\star\star$ \iqroles{trader}]\label{iq:m3:metals:5}
Pourquoi l'EFP sur l'or s'est-il creusé en mars 2020, et qu'est-ce qui l'a refermé~?
\end{interviewq}

\begin{interviewq}[$\star\star\star$ \iqroles{développeur, risque}]\label{iq:m3:metals:6}
Vous supervisez la surveillance d'une bourse. Quels signaux auraient dû vous alerter dans la nuit du 7 au 8 mars
2022, et qu'auriez-vous automatisé~?
\end{interviewq}
